% Автоматаар үүсгэсэн — эх файл: lab08/k3s/README.md
\chapter{K3s манифестууд}

XIV долоо хоногийн бие даалтаар эдгээрийг \textbf{УНШИЖ ойлгосон} байх ёстой.

\hypertarget{ux44eux443ux433-ux445ux430ux430ux43dux430-ux430ux436ux438ux43bux43bux443ux443ux43bux430ux445-ux432ux44d}{%
\section{Юуг хаана ажиллуулах вэ}\label{ux44eux443ux433-ux445ux430ux430ux43dux430-ux430ux436ux438ux43bux43bux443ux443ux43bux430ux445-ux432ux44d}}

\begingroup\scriptsize\begin{longtable}[]{@{}
  >{\raggedright\arraybackslash}p{(\columnwidth - 4\tabcolsep) * \real{0.3333}}
  >{\raggedright\arraybackslash}p{(\columnwidth - 4\tabcolsep) * \real{0.3333}}
  >{\raggedright\arraybackslash}p{(\columnwidth - 4\tabcolsep) * \real{0.3333}}@{}}
\toprule
\begin{minipage}[b]{\linewidth}\raggedright
Давхарга
\end{minipage} & \begin{minipage}[b]{\linewidth}\raggedright
Төмөр
\end{minipage} & \begin{minipage}[b]{\linewidth}\raggedright
Юу ажиллах вэ
\end{minipage} \\
\midrule
\endhead
\bottomrule
\endlastfoot
Ирмэг & \textbf{Raspberry Pi 3B} (K3s) & mosquitto (гүүр) + edge-agent \\
Үүл & зөөврийн компьютер (Docker Compose) & EMQX, InfluxDB, Grafana, registry, Node-RED, Dex, GraphQL, Ollama \\
\end{longtable}\endgroup

\begin{quote}
\textbf{EMQX/InfluxDB/Grafana-г Pi-гийн K3s рүү оруулах гэж бүү оролд.} Зөвхөн
InfluxDB 3 нь 250--500 MiB, EMQX 150--250 MiB иддэг. K3s өөрөө \textasciitilde400 MiB
авсны дараа 1 GB-д тэдгээр багтахгүй --- pod нь \texttt{OOMKilled} эсвэл
\texttt{Pending} болно. Энэ нь алдаа биш, \textbf{тоо ярьж байна}.
\end{quote}

\hypertarget{ux444ux430ux439ux43bux443ux443ux434}{%
\section{Файлууд}\label{ux444ux430ux439ux43bux443ux443ux434}}

\begingroup\scriptsize\begin{longtable}[]{@{}
  >{\raggedright\arraybackslash}p{(\columnwidth - 2\tabcolsep) * \real{0.5000}}
  >{\raggedright\arraybackslash}p{(\columnwidth - 2\tabcolsep) * \real{0.5000}}@{}}
\toprule
\begin{minipage}[b]{\linewidth}\raggedright
Файл
\end{minipage} & \begin{minipage}[b]{\linewidth}\raggedright
Агуулга
\end{minipage} \\
\midrule
\endhead
\bottomrule
\endlastfoot
\texttt{00-namespace.yaml} & \texttt{cnc302} нэрийн орон зай \\
\texttt{01-config.yaml} & ConfigMap (UNS + \texttt{CLOUD\_HOST}) + Secret загвар \\
\texttt{10-mosquitto.yaml} & ConfigMap (mosquitto.conf + гүүрний загвар) + PVC + Deployment + NodePort Service \\
\texttt{20-edge-agent.yaml} & Deployment (локал импортолсон дүрс) \\
\texttt{90-hpa.yaml} & HPA --- ба яагаад ганц зангилаа дээр 1→2-оос цааш явахгүй нь \\
\end{longtable}\endgroup

\hypertarget{ux441ux430ux43dux430ux445-ux43eux439ux43d-ux442ux43eux43eux446ux43eux43e-ux44dux43dux44d-ux445ux44dux441ux433ux438ux439ux433-ux430ux43bux433ux430ux441ux430ux436-ux431ux43eux43bux43eux445ux433ux4afux439}{%
\section{1. Санах ойн тооцоо (ЭНЭ ХЭСГИЙГ АЛГАСАЖ БОЛОХГҮЙ)}\label{ux441ux430ux43dux430ux445-ux43eux439ux43d-ux442ux43eux43eux446ux43eux43e-ux44dux43dux44d-ux445ux44dux441ux433ux438ux439ux433-ux430ux43bux433ux430ux441ux430ux436-ux431ux43eux43bux43eux445ux433ux4afux439}}

\begin{verbatim}
Pi 3B-гийн нийт RAM                            1024 MiB
  − GPU-д хуваарилсан (gpu_mem=64)              −64 MiB
  − firmware / kernel-ийн эзэмшил               −35 MiB
  ────────────────────────────────────────────────────
  MemTotal (free -m харуулах утга)            ≈ 925 MiB
  − Raspberry Pi OS Lite (systemd, sshd, …)     −45 MiB
  ────────────────────────────────────────────────────
  БОЛОМЖТОЙ                                   ≈ 880 MiB
\end{verbatim}

Үүнээс:

\begingroup\scriptsize\begin{longtable}[]{@{}
  >{\raggedright\arraybackslash}p{(\columnwidth - 4\tabcolsep) * \real{0.3000}}
  >{\raggedleft\arraybackslash}p{(\columnwidth - 4\tabcolsep) * \real{0.4000}}
  >{\raggedright\arraybackslash}p{(\columnwidth - 4\tabcolsep) * \real{0.3000}}@{}}
\toprule
\begin{minipage}[b]{\linewidth}\raggedright
Хэрэглэгч
\end{minipage} & \begin{minipage}[b]{\linewidth}\raggedleft
Бодит RSS
\end{minipage} & \begin{minipage}[b]{\linewidth}\raggedright
Тэмдэглэл
\end{minipage} \\
\midrule
\endhead
\bottomrule
\endlastfoot
k3s server (apiserver + kubelet + containerd + sqlite) & \textbf{\textasciitilde400 MiB} & pod биш, systemd үйлчилгээ --- \texttt{requests}-д ОРОХГҮЙ \\
coredns & \textasciitilde25 MiB & request 70Mi (k3s-ийн анхны утга) \\
metrics-server & \textasciitilde25 MiB & HPA-д ЗААВАЛ хэрэгтэй \\
local-path-provisioner & \textasciitilde10 MiB & PVC-г microSD дээр үүсгэнэ \\
\textbf{mosquitto} (манай) & \textasciitilde12 MiB & request 32Mi / limit 96Mi \\
\textbf{edge-agent} (манай) & \textasciitilde55--110 MiB & request 96Mi / limit 160Mi \\
\textbf{Нийт (1 агент)} & \textbf{≈ 530 MiB} & \textasciitilde350 MiB нөөц үлдэнэ \\
\textbf{Нийт (HPA 2 агент)} & \textbf{≈ 600 MiB} & \textasciitilde280 MiB нөөц --- ЭНЭ Л ХЯЗГААР \\
\end{longtable}\endgroup

Товлогчийн (scheduler) тал дээрх тооцоо өөр байдгийг анзаар:

\begin{verbatim}
allocatable ≈ MemTotal − eviction threshold (100Mi) ≈ 825 MiB
хүсэлтүүд (requests): coredns 70 + metrics-server 70 + mosquitto 32
                      + edge-agent 96 × N
  N = 1 → 268 Mi     N = 2 → 364 Mi     N = 3 → 460 Mi
\end{verbatim}

\texttt{requests}-ийн хувьд 3 хувь ч "багтана", \textbf{гэвч бодит RSS багтахгүй} ---
яг үүнээс болж 3 дахь pod нь Pending биш, \texttt{OOMKilled} болж эхэлдэг.
\texttt{requests} бол амлалт, \texttt{limits} бол хана, \textbf{RSS бол үнэн}. Гурвыг нь
\texttt{kubectl\ top} + \texttt{free\ -m}-ээр зэрэг хэмжиж тайландаа бич.

\hypertarget{k3s-ux441ux443ux443ux43bux433ux430ux445-ux43dux44dux433-ux443ux434ux430ux430-pi-3b-ux434ux44dux44dux440}{%
\section{2. K3s суулгах (нэг удаа, Pi 3B дээр)}\label{k3s-ux441ux443ux443ux43bux433ux430ux445-ux43dux44dux433-ux443ux434ux430ux430-pi-3b-ux434ux44dux44dux440}}

\textbf{Эхлээд Pi дээрх Docker Compose-ыг БҮРЭН зогсооно} --- эс тэгвээс хоёулаа
нэг санах ойг булаацалдана:

\begin{Shaded}
\begin{Highlighting}[]
\BuiltInTok{cd}\NormalTok{ \textasciitilde{}/cnc302/edge }\KeywordTok{\&\&} \ExtensionTok{docker}\NormalTok{ compose down}
\ExtensionTok{docker}\NormalTok{ ps                      }\CommentTok{\# хоосон байх ёстой}
\FunctionTok{free} \AttributeTok{{-}m}                        \CommentTok{\# 800+ MiB сул байх ёстой}
\end{Highlighting}
\end{Shaded}

Swap-г 1 GB болгож нэмбэл ачаалалтай үед аврах боловч microSD-г элээнэ:

\begin{Shaded}
\begin{Highlighting}[]
\FunctionTok{sudo}\NormalTok{ dphys{-}swapfile swapoff}
\FunctionTok{sudo}\NormalTok{ sed }\AttributeTok{{-}i} \StringTok{\textquotesingle{}s/\^{}CONF\_SWAPSIZE=.*/CONF\_SWAPSIZE=1024/\textquotesingle{}}\NormalTok{ /etc/dphys{-}swapfile}
\FunctionTok{sudo}\NormalTok{ dphys{-}swapfile setup }\KeywordTok{\&\&} \FunctionTok{sudo}\NormalTok{ dphys{-}swapfile swapon}
\end{Highlighting}
\end{Shaded}

Суулгалт:

\begin{Shaded}
\begin{Highlighting}[]
\ExtensionTok{curl} \AttributeTok{{-}sfL}\NormalTok{ https://get.k3s.io }\KeywordTok{|} \FunctionTok{sh} \AttributeTok{{-}s} \AttributeTok{{-}} \DataTypeTok{\textbackslash{}}
  \AttributeTok{{-}{-}write{-}kubeconfig{-}mode}\NormalTok{ 644 }\DataTypeTok{\textbackslash{}}
  \AttributeTok{{-}{-}disable}\NormalTok{ traefik }\DataTypeTok{\textbackslash{}}
  \AttributeTok{{-}{-}disable}\NormalTok{ servicelb}
\end{Highlighting}
\end{Shaded}

\begin{itemize}
\tightlist
\item
  \texttt{-\/-disable\ traefik} → Ingress контроллер \textasciitilde60 MiB иднэ, бидэнд MQTT
  (TCP) хэрэгтэй болохоос HTTP Ingress хэрэггүй.
\item
  \texttt{-\/-disable\ servicelb} → LoadBalancer төрлийн Service ашиглахгүй,
  NodePort хангалттай (\textasciitilde15 MiB хэмнэнэ).
\item
  \textbf{metrics-server-ийг УНТРААХГҮЙ} --- HPA түүнгүйгээр ажиллахгүй.
\end{itemize}

\begin{Shaded}
\begin{Highlighting}[]
\BuiltInTok{export} \VariableTok{KUBECONFIG}\OperatorTok{=}\NormalTok{/etc/rancher/k3s/k3s.yaml}
\ExtensionTok{kubectl}\NormalTok{ get nodes                 }\CommentTok{\# Ready болтол 1–3 минут (microSD удаан)}
\ExtensionTok{kubectl}\NormalTok{ top nodes                 }\CommentTok{\# metrics{-}server бэлэн эсэх}
\end{Highlighting}
\end{Shaded}

\hypertarget{ux434ux4afux440ux441ux438ux439ux433-pi-ux434ux44dux44dux440-ux431ux430ux440ux44cux436-ux438ux43cux43fux43eux440ux442ux43bux43eux445}{%
\section{3. Дүрсийг Pi дээр барьж импортлох}\label{ux434ux4afux440ux441ux438ux439ux433-pi-ux434ux44dux44dux440-ux431ux430ux440ux44cux436-ux438ux43cux43fux43eux440ux442ux43bux43eux445}}

K3s нь Docker-ын дүрсийн санг ХАРАХГҮЙ (өөрийн containerd-тэй). Тиймээс:

\begin{Shaded}
\begin{Highlighting}[]
\BuiltInTok{cd}\NormalTok{ \textasciitilde{}/cnc302/edge/agent}
\ExtensionTok{docker}\NormalTok{ build }\AttributeTok{{-}t}\NormalTok{ cnc302/edge{-}agent:v1 .              }\CommentTok{\# Pi 3B дээр 3–6 минут}
\ExtensionTok{docker}\NormalTok{ save cnc302/edge{-}agent:v1 }\KeywordTok{|} \FunctionTok{sudo}\NormalTok{ k3s ctr images import }\AttributeTok{{-}}
\FunctionTok{sudo}\NormalTok{ k3s ctr images ls }\KeywordTok{|} \FunctionTok{grep}\NormalTok{ edge{-}agent            }\CommentTok{\# байгаа эсэхийг шалга}
\end{Highlighting}
\end{Shaded}

\texttt{docker} байхгүй бол шууд tar-аар:

\begin{Shaded}
\begin{Highlighting}[]
\FunctionTok{sudo}\NormalTok{ k3s ctr images import cnc302{-}edge{-}agent{-}v1.tar}
\end{Highlighting}
\end{Shaded}

Мөн \texttt{busybox:1.36} (initContainer) хэрэгтэй. Интернэт байвал автоматаар
татагдана, үгүй бол урьдчилан импортлоно:

\begin{Shaded}
\begin{Highlighting}[]
\FunctionTok{sudo}\NormalTok{ k3s ctr images pull docker.io/library/busybox:1.36}
\FunctionTok{sudo}\NormalTok{ k3s ctr images pull docker.io/library/eclipse{-}mosquitto:2.0.20}
\end{Highlighting}
\end{Shaded}

\begin{quote}
\texttt{imagePullPolicy:\allowbreak{}\ IfNotPresent} нь энэ бүх ажлын үндэс. \texttt{Always} болговол
K3s Docker Hub-аас \texttt{cnc302/edge-agent} хайж олохгүй → \texttt{ErrImagePull}.
\end{quote}

\hypertarget{ux431ux430ux439ux440ux448ux443ux443ux43bux430ux445}{%
\section{4. Байршуулах}\label{ux431ux430ux439ux440ux448ux443ux443ux43bux430ux445}}

\begin{Shaded}
\begin{Highlighting}[]
\BuiltInTok{cd}\NormalTok{ \textasciitilde{}/cnc302/lab08/k3s}
\CommentTok{\# ⚠ 01{-}config.yaml дотор CLOUD\_HOST{-}ыг компьютерийнхээ LAN IP болгож солино!}
\ExtensionTok{kubectl}\NormalTok{ apply }\AttributeTok{{-}f}\NormalTok{ .}

\ExtensionTok{kubectl} \AttributeTok{{-}n}\NormalTok{ cnc302 get pods }\AttributeTok{{-}w}
\ExtensionTok{kubectl} \AttributeTok{{-}n}\NormalTok{ cnc302 logs deploy/mosquitto }\AttributeTok{{-}c}\NormalTok{ render{-}bridge   }\CommentTok{\# гүүрний тохиргоо}
\ExtensionTok{kubectl} \AttributeTok{{-}n}\NormalTok{ cnc302 logs }\AttributeTok{{-}f}\NormalTok{ deploy/edge{-}agent}
\end{Highlighting}
\end{Shaded}

Ажиллаж байгаа эсэхийг \textbf{үүлний талаас} батал:

\begin{Shaded}
\begin{Highlighting}[]
\CommentTok{\# зөөврийн компьютер дээр}
\ExtensionTok{mosquitto\_sub} \AttributeTok{{-}h}\NormalTok{ localhost }\AttributeTok{{-}t} \StringTok{\textquotesingle{}cnc302/shutis/\#\textquotesingle{}} \AttributeTok{{-}v} \AttributeTok{{-}C}\NormalTok{ 10}
\end{Highlighting}
\end{Shaded}

Гүүрний төлөв:

\begin{Shaded}
\begin{Highlighting}[]
\ExtensionTok{mosquitto\_sub} \AttributeTok{{-}h}\NormalTok{ localhost }\AttributeTok{{-}t} \StringTok{\textquotesingle{}cnc302/shutis/mhts/lab/pi3b{-}01/bridge/state\textquotesingle{}} \AttributeTok{{-}v}
\CommentTok{\# 1 = холбогдсон, 0 = тасарсан}
\end{Highlighting}
\end{Shaded}

\hypertarget{ux445ux44dux43cux436ux438ux445}{%
\section{5. Хэмжих}\label{ux445ux44dux43cux436ux438ux445}}

\begin{Shaded}
\begin{Highlighting}[]
\ExtensionTok{kubectl}\NormalTok{ top nodes}
\ExtensionTok{kubectl}\NormalTok{ top pods }\AttributeTok{{-}n}\NormalTok{ cnc302}
\FunctionTok{free} \AttributeTok{{-}m}                                   \CommentTok{\# kubectl top{-}той харьцуул}
\ExtensionTok{kubectl} \AttributeTok{{-}n}\NormalTok{ cnc302 describe node }\KeywordTok{|} \FunctionTok{sed} \AttributeTok{{-}n} \StringTok{\textquotesingle{}/Allocated resources/,/\^{}Events/p\textquotesingle{}}
\end{Highlighting}
\end{Shaded}

\textbf{Хүснэгт (тайланд):} үйлчилгээ бүрийн \texttt{requests} / \texttt{limits} / бодит RSS,
мөн Compose дээрх ижил үйлчилгээний \texttt{docker\ stats} утга. Ялгааг тайлбарла.

HPA-г ажиглах:

\begin{Shaded}
\begin{Highlighting}[]
\ExtensionTok{kubectl} \AttributeTok{{-}n}\NormalTok{ cnc302 get hpa edge{-}agent }\AttributeTok{{-}w}
\CommentTok{\# өөр терминалд ачаалал үүсгэ:}
\ExtensionTok{kubectl} \AttributeTok{{-}n}\NormalTok{ cnc302 exec deploy/edge{-}agent }\AttributeTok{{-}{-}}\NormalTok{ python }\AttributeTok{{-}c} \StringTok{"}
\StringTok{while True: pass"} \KeywordTok{\&}
\end{Highlighting}
\end{Shaded}

\hypertarget{pod-ux43dux44c-oomkilled-ux431ux43eux43bux432ux43eux43b-ux44eux443-ux445ux438ux439ux445-ux432ux44d}{%
\section{6. Pod нь OOMKilled болвол юу хийх вэ}\label{pod-ux43dux44c-oomkilled-ux431ux43eux43bux432ux43eux43b-ux44eux443-ux445ux438ux439ux445-ux432ux44d}}

Эхлээд \textbf{баримтыг цуглуул} (тайланд хэрэгтэй):

\begin{Shaded}
\begin{Highlighting}[]
\ExtensionTok{kubectl} \AttributeTok{{-}n}\NormalTok{ cnc302 get pods                       }\CommentTok{\# RESTARTS багана}
\ExtensionTok{kubectl} \AttributeTok{{-}n}\NormalTok{ cnc302 describe pod }\OperatorTok{\textless{}}\NormalTok{pod}\OperatorTok{\textgreater{}} \KeywordTok{|} \FunctionTok{grep} \AttributeTok{{-}A3} \StringTok{"Last State"}
\CommentTok{\#   Last State: Terminated,  Reason: OOMKilled,  Exit Code: 137}
\ExtensionTok{kubectl}\NormalTok{ get events }\AttributeTok{{-}n}\NormalTok{ cnc302 }\AttributeTok{{-}{-}sort{-}by}\OperatorTok{=}\NormalTok{.lastTimestamp }\KeywordTok{|} \FunctionTok{tail} \AttributeTok{{-}20}
\FunctionTok{dmesg} \AttributeTok{{-}T} \KeywordTok{|} \FunctionTok{grep} \AttributeTok{{-}i} \StringTok{"killed process"}              \CommentTok{\# kernel{-}ийн OOM killer}
\end{Highlighting}
\end{Shaded}

Дараа нь \textbf{шалтгаанаар нь} ялга:

\begingroup\scriptsize\begin{longtable}[]{@{}
  >{\raggedright\arraybackslash}p{(\columnwidth - 4\tabcolsep) * \real{0.3333}}
  >{\raggedright\arraybackslash}p{(\columnwidth - 4\tabcolsep) * \real{0.3333}}
  >{\raggedright\arraybackslash}p{(\columnwidth - 4\tabcolsep) * \real{0.3333}}@{}}
\toprule
\begin{minipage}[b]{\linewidth}\raggedright
Шинж
\end{minipage} & \begin{minipage}[b]{\linewidth}\raggedright
Шалтгаан
\end{minipage} & \begin{minipage}[b]{\linewidth}\raggedright
Хийх зүйл
\end{minipage} \\
\midrule
\endhead
\bottomrule
\endlastfoot
Exit 137, \texttt{Reason:\ OOMKilled} & контейнер өөрийн \texttt{limits}-ээ давсан & \texttt{limits.memory}-г 32--64Mi-аар нэм, ЭСВЭЛ ажлыг хөнгөл (жишээ: \texttt{INFER\_THREADS=1}, \texttt{INTERVAL}-ыг өсгө) \\
Pod \texttt{Pending}, \texttt{Insufficient\ memory} & зангилаанд \texttt{requests} багтахгүй & хувийн тоог бууруул, эсвэл \texttt{requests}-ээ бодит RSS-д ойртуул \\
Зангилаа өөрөө \texttt{NotReady}, k3s унтарсан & системийн OOM killer k3s-ийг алсан & Compose зогссон эсэхийг шалга, swap нэм, \texttt{-\/-disable}-уудыг батал \\
Байнгын CrashLoopBackOff & дүрс буруу архитектур / импортлоогүй & \texttt{sudo\ k3s\ ctr\ images\ ls}, дүрсийг Pi дээр дахин бари \\
\end{longtable}\endgroup

\textbf{Гол зарчим:} \texttt{limits}-ийг ӨСГӨХ нь үргэлж зөв шийдэл БИШ. 1 GB дээр
хамгийн зөв хариулт нь ихэвчлэн "энэ ажлыг Pi дээр биш, үүлэн дээр
ажиллуул" байдаг --- Лаб 8-ын гол сургамж яг энэ.

\hypertarget{compose-ux431ux430-kubernetes-ux438ux439ux43d-ux445ux430ux440ux44cux446ux443ux443ux43bux430ux43bux442}{%
\section{7. Compose ба Kubernetes-ийн харьцуулалт}\label{compose-ux431ux430-kubernetes-ux438ux439ux43d-ux445ux430ux440ux44cux446ux443ux443ux43bux430ux43bux442}}

\begingroup\scriptsize\begin{longtable}[]{@{}
  >{\raggedright\arraybackslash}p{(\columnwidth - 4\tabcolsep) * \real{0.3333}}
  >{\raggedright\arraybackslash}p{(\columnwidth - 4\tabcolsep) * \real{0.3333}}
  >{\raggedright\arraybackslash}p{(\columnwidth - 4\tabcolsep) * \real{0.3333}}@{}}
\toprule
\begin{minipage}[b]{\linewidth}\raggedright
Ойлголт
\end{minipage} & \begin{minipage}[b]{\linewidth}\raggedright
Docker Compose (\texttt{edge/})
\end{minipage} & \begin{minipage}[b]{\linewidth}\raggedright
Kubernetes (энэ фолдер)
\end{minipage} \\
\midrule
\endhead
\bottomrule
\endlastfoot
Үйлчилгээ & \texttt{services:} & Deployment + Service \\
Боть & \texttt{volumes:} & PersistentVolumeClaim (\texttt{local-path}) \\
Орчны хувьсагч & \texttt{.env} + \texttt{environment:} & ConfigMap / Secret \\
Тохиргооны файл & bind mount & ConfigMap + \texttt{subPath} \\
Загвар орлуулах & \texttt{make\ bridge} (host дээр \texttt{sed}) & \texttt{initContainer} + \texttt{sed} \\
Эрүүл мэнд & \texttt{healthcheck:} & readinessProbe / livenessProbe \\
Нөөцийн хязгаар & \texttt{deploy.resources.limits} & \texttt{resources.requests} \textbf{ба} \texttt{limits} \\
Дахин эхлүүлэх & \texttt{restart:} & ReplicaSet (автоматаар) \\
Порт нээх & \texttt{ports:} & NodePort \\
Өргөтгөх & гараар \texttt{-\/-scale} & \texttt{replicas} эсвэл HPA \\
\end{longtable}\endgroup

\textbf{\texttt{requests} ба \texttt{limits}-ийн ялгаа чухал:} \texttt{requests} нь товлогчид "энэ
pod-д хамгийн багадаа ийм нөөц хэрэгтэй" гэж хэлнэ (товлолт үүгээр
шийдэгдэнэ); \texttt{limits} нь дээд хана (давбал OOMKilled). Compose-д зөвхөн
хана байдаг --- тиймээс Compose нь "багтах уу, үгүй юу"-г УРЬДЧИЛАН хэлж
чадахгүй, зөвхөн унасны дараа мэдэгддэг.

\hypertarget{ux443ux441ux442ux433ux430ux445-ux431ux443ux446ux430ux430ux445}{%
\section{8. Устгах / буцаах}\label{ux443ux441ux442ux433ux430ux445-ux431ux443ux446ux430ux430ux445}}

\begin{Shaded}
\begin{Highlighting}[]
\ExtensionTok{kubectl}\NormalTok{ delete }\AttributeTok{{-}f}\NormalTok{ .}
\ExtensionTok{kubectl}\NormalTok{ delete namespace cnc302        }\CommentTok{\# PVC{-}тэй хамт — өгөгдөл алдагдана!}

\CommentTok{\# K3s{-}ийг бүрэн устгаж Compose руу буцах}
\FunctionTok{sudo}\NormalTok{ /usr/local/bin/k3s{-}uninstall.sh}
\BuiltInTok{cd}\NormalTok{ \textasciitilde{}/cnc302/edge }\KeywordTok{\&\&} \FunctionTok{make}\NormalTok{ up}
\end{Highlighting}
\end{Shaded}
